\documentclass[11pt]{article}
\usepackage[margin=1in]{geometry}
\usepackage{xcolor}
\usepackage{tcolorbox}
\usepackage{hyperref}
\usepackage{enumitem}
\setlist{noitemsep,topsep=0pt}
\usepackage{booktabs}
\usepackage{array}
\usepackage{amsmath}

% Define OSU orange color
\definecolor{osaorange}{RGB}{220,115,38}

% Configure hyperref colors
\hypersetup{
    colorlinks=true,
    linkcolor=blue,
    urlcolor=blue,
    citecolor=blue
}

\title{}
\author{}
\date{}

\begin{document}

% Orange header box with black outline
\begin{tcolorbox}[colback=osaorange,colframe=black,boxrule=2pt,arc=0pt,left=6pt,right=6pt,top=10pt,bottom=10pt]
\centering
\textcolor{white}{\textbf{\Large Week 5 In-Class Worksheet}}\\
\textcolor{white}{\textbf{\large Hypothesis Testing}}
\end{tcolorbox}

\vspace{8pt}

\noindent\textbf{Name:} \rule{8cm}{0.4pt} \hfill \textbf{Date:} \rule{4cm}{0.4pt}\\[0.5ex]
\textbf{Course:} H524 Introduction to Biostatistics\\
\textbf{Topic:} Hypothesis Testing for Population Means

\vspace{8pt}

\section*{Instructions}
Work with a partner to solve the following problems. Show all your work, including formulas and calculations. You may use R to verify your answers, but demonstrate your understanding by showing the steps manually first.

\vspace{12pt}

\section{Problem 1: Two-Sided t-test}

A researcher claims that the average body temperature is not 98.6°F. A sample of 25 healthy adults has a mean temperature of 98.4°F with a standard deviation of 0.8°F. Test the claim at the 0.05 significance level.

\textbf{Given:} $n = 25$, $\bar{x} = 98.4$°F, $s = 0.8$°F, $\alpha = 0.05$

\vspace{6pt}

\textbf{Part A:} State the null and alternative hypotheses.

\vspace{3cm}

\textbf{Part B:} Calculate the standard error and t-statistic.

\vspace{0.3cm}

\textbf{Formula:} $\text{SE} = \frac{s}{\sqrt{n}}$

\vspace{2cm}

\textbf{Formula:} $t = \frac{\bar{x} - \mu_0}{\text{SE}}$

\vspace{2cm}

\textbf{Part C:} Find the p-value.

\vspace{0.3cm}

Degrees of freedom: $\text{df} = n - 1 = $

\vspace{0.5cm}

For a two-sided test, p-value = $2 \times P(T < t)$ where $t$ is your test statistic.

(Use R: \texttt{2 * pt(t, df=24)})

\vspace{2cm}

\textbf{Part D:} Make a decision and state your conclusion in context.

\vspace{3cm}

\newpage

\section{Problem 2: One-Sided t-test (Upper Tail)}

A nutritionist claims that a new diet increases average daily caloric expenditure above the baseline of 260 calories. A study of 40 participants following the diet shows a mean expenditure of 285 calories with a standard deviation of 48 calories. Test this claim at the 0.05 level.

\textbf{Given:} $n = 40$, $\bar{x} = 285$, $s = 48$, $\mu_0 = 260$, $\alpha = 0.05$

\vspace{6pt}

\textbf{Part A:} State the hypotheses and identify the type of test.

\vspace{3cm}

\textbf{Part B:} Calculate the standard error and t-statistic.

\vspace{0.3cm}

Standard error: $\text{SE} = \frac{s}{\sqrt{n}} = $

\vspace{2cm}

t-statistic: $t = \frac{\bar{x} - \mu_0}{\text{SE}} = $

\vspace{2cm}

\textbf{Part C:} Find the p-value and make a decision.

\vspace{0.3cm}

Degrees of freedom: $\text{df} = $

\vspace{0.5cm}

For a one-sided (upper tail) test: p-value = $P(T > t)$

(Use R: \texttt{1 - pt(t, df)})

\vspace{2cm}

Decision:

\vspace{1.5cm}

\textbf{Part D:} Interpret the results in context. What does this tell you about the diet's effect?

\vspace{3cm}

\newpage

\section{Problem 3: Critical Value Approach}

Using the data from Problem 2, verify your conclusion using the critical value approach instead of the p-value approach.

\textbf{Data:} $n=40$, $\bar{x}=285$, $s=48$, $t=$[your calculated value], $\alpha=0.05$

\vspace{6pt}

\textbf{Part A:} Find the critical value for this one-sided test.

\vspace{0.3cm}

For a one-sided (upper tail) test with $\alpha = 0.05$ and df = 39:

(Use R: \texttt{qt(0.95, df=39)})

\vspace{2cm}

\textbf{Part B:} Compare the t-statistic to the critical value and make a decision.

\vspace{3cm}

\textbf{Part C:} Compare this conclusion to the p-value approach from Problem 2. Do they agree?

\vspace{3cm}

\newpage

\section{Problem 4: Type I and Type II Errors}

A pharmaceutical company is testing a new treatment that they hope is more effective than the standard treatment.

\vspace{6pt}

\textbf{Part A:} Define a Type I error in the context of this study. What would be the consequence?

\vspace{4cm}

\textbf{Part B:} Define a Type II error in the context of this study. What would be the consequence?

\vspace{4cm}

\textbf{Part C:} Which type of error do you think is more serious in medical research? Explain your reasoning.

\vspace{4cm}

\newpage

\section{Problem 5: Statistical Power and Sample Size}

Answer the following conceptual questions about statistical power.

\vspace{6pt}

\textbf{Part A:} How does increasing the significance level (e.g., from 0.05 to 0.10) affect the probability of Type I and Type II errors? What happens to statistical power?

\vspace{5cm}

\textbf{Part B:} How does increasing sample size affect Type I error, Type II error, and statistical power? Explain the mechanism.

\vspace{5cm}

\newpage

\section{Problem 6: Comprehensive Integration Problem}

A pharmaceutical company is testing whether a new sleep medication reduces the average sleep time below the standard 8 hours. They conduct a study with 50 patients and find a mean sleep time of 7.8 hours with a standard deviation of 1.2 hours. Test the hypothesis at the $\alpha = 0.01$ level.

\textbf{Given:} $n = 50$, $\bar{x} = 7.8$ hours, $s = 1.2$ hours, $\alpha = 0.01$

\vspace{6pt}

\textbf{Part A:} State the null and alternative hypotheses. Is this a one-sided or two-sided test?

\vspace{3cm}

\textbf{Part B:} Calculate the standard error.

\vspace{0.3cm}

\textbf{Formula:} $\text{SE} = \frac{s}{\sqrt{n}}$

\vspace{2cm}

\textbf{Part C:} Calculate the t-statistic.

\vspace{0.3cm}

\textbf{Formula:} $t = \frac{\bar{x} - \mu_0}{\text{SE}}$

\vspace{2cm}

\textbf{Part D:} Find the p-value.

\vspace{0.3cm}

Degrees of freedom: $\text{df} = $

\vspace{0.5cm}

For a one-sided (lower tail) test: p-value = $P(T < t)$

(Use R: \texttt{pt(t, df)})

\vspace{2cm}

\textbf{Part E:} Make a decision at the $\alpha = 0.01$ level.

\vspace{2cm}

\textbf{Part F:} Write a complete interpretation of your results in the context of the sleep study.

\vspace{3cm}

\textbf{Part G:} Verify your conclusion using the critical value approach.

\vspace{0.3cm}

For a one-sided (lower tail) test with $\alpha = 0.01$ and df = 49:

(Use R: \texttt{qt(0.01, df=49)})

\vspace{1cm}

Compare your t-statistic to the critical value:

\vspace{3cm}

\newpage

\section*{Reflection Questions}

\textbf{1.} Explain the difference between statistical significance and practical significance. Can a result be statistically significant but not practically important?

\vspace{4cm}

\textbf{2.} Why do we typically use $\alpha = 0.05$ as the significance level? What would happen if we used $\alpha = 0.10$ or $\alpha = 0.01$ instead?

\vspace{4cm}

\textbf{3.} When would you use a one-sided test versus a two-sided test? Give an example of each.

\vspace{4cm}

\textbf{4.} What is the relationship between confidence intervals and hypothesis testing?

\vspace{4cm}

\end{document}
